\section{The first quotient and the second-layer different}\label{sec:ramification}

We now calculate the actual ramification invariants. The
element \(\alpha_2-\beta_1\) has valuation divisible by \(p\)
in \(L_2L_1\), so Lemma~\ref{lem:prime-value} does not apply
there. Its replacement tracks the leading cancellation.

\begin{lemma}[A once-\(p\) cancellation]\label{lem:cancellation}
Let \(D/K\) be a finite totally ramified Galois \(p\)-extension,
with group \(G\), first lower break \(b_0>0\), and \(p\nmid b_0\).
Suppose \(x\in D\) has \(a=v_D(x)>0\), \(v_p(a)=1\), and
\begin{equation}\label{eq:cancellation-hyp}
 v_D(gx-x)>a+p b_0\qquad(g\ne1).
\end{equation}
In any uniformizer expansion of \(x\), the first nonzero
exponent prime to \(p\) is
\[
 n=a+(p-1)b_0.
\]
The first graded quotient \(G_{b_0}/G_{b_0+1}\) has order \(p\).
If \(g\) has later lower break \(c=\ell_D(g)>b_0\), then
\begin{equation}\label{eq:later-rule}
 v_D(gx-x)=n+c.
\end{equation}
\end{lemma}

\begin{proof}
Write \(x=\sum_{j\ge a}c_j\pi^j\), with \(c_a\ne0\).
For an element \(g\) of break \(b_0\),
\eqref{eq:monomial-displacement} says that the leading
monomial contributes at exact order \(a+p b_0\).
Every later monomial with exponent divisible by \(p\)
contributes at strictly greater order.
If the first nonzero exponent \(j\) prime to \(p\) were
less than \(a+(p-1)b_0\), its contribution at \(j+b_0\)
would be uniquely lowest, contradicting
\eqref{eq:cancellation-hyp}.
If no such exponent were present at \(n=a+(p-1)b_0\),
the leading monomial at order \(a+p b_0\) could not cancel.
Thus \(n\) is the first such exponent, and \(p\nmid n\).

To control the entire first grade, not just one automorphism, put
\[
 \theta(g)=\res\left(\frac{g\pi-\pi}{\pi^{b_0+1}}\right),
 \qquad g\in G=G_{b_0}.
\]
Composition adds the leading coefficients, since each
automorphism fixes \(k\) and has linear coefficient \(1\).
Hence \(\theta\) is additive and has kernel \(G_{b_0+1}\).
The coefficient of order \(a+p b_0\) in \(gx-x\) is
\[
 c_a\,\overline{(a/p)}\,\theta(g)^p
 +c_n\overline n\,\theta(g).
\]
Only these two monomials contribute at that order.
By \eqref{eq:cancellation-hyp}, the displayed expression
vanishes for every \(g\). Its leading coefficient is nonzero,
so this degree-\(p\) polynomial has at most \(p\) roots.
The image of the nontrivial first graded \(p\)-group therefore
has exactly \(p\) elements.

For a later break \(c>b_0\), the monomial \(c_n\pi^n\)
contributes at order \(n+c\). Later prime-to-\(p\) exponents
give larger orders. Every exponent divisible by \(p\)
gives order at least \(a+pc\), and
\[
 a+pc-(n+c)=(p-1)(c-b_0)>0.
\]
The contribution at \(n+c\) is unique. All tail orders tend
to infinity, so no infinite-series cancellation changes it.
This proves \eqref{eq:later-rule}.
\end{proof}

\subsection{Apply the cancellation in the disjoint compositum}

Set \(B=L_2L_1\), \(G=\Gal(B/K)\), and choose
\(\alpha\in S_2,\beta\in S_1\). By
Proposition~\ref{prop:separation},
\[
 G\simeq C_{p^2}\times C_p,\qquad v_B=p^3v.
\]
Let \(\sigma,\gamma\in G\) act by one native step on
\(\alpha,\beta\), respectively, fixing the other root.
For \(\eta=\alpha-\beta\), write
\begin{equation}\label{eq:a-A}
 a=v_B(\eta)=2p(p-1)^2,\qquad A=2p^2(p-1).
\end{equation}
The contact gives the first equality, with \(v_p(a)=1\).
Both root displacements have base valuation at least \(2r\),
so
\begin{equation}\label{eq:uniform-displ}
 v_B(g\eta-\eta)\ge A\quad(g\ne1),\qquad
 v_B(\sigma\eta-\eta)=v_B(\gamma\eta-\eta)=A.
\end{equation}
Infinite valuation is allowed in the inequality.

Let \(b_0\) be the first lower, equivalently upper, break
of \(B/K\). Its quotient \(L_1/K\) gives \(b_0\le p-1\).
The first drop of a finite abelian \(p\)-extension is detected
by a degree-\(p\) character: the first graded quotient is
elementary abelian (the leading-coefficient injection above
shows this), and a nonzero character of that quotient pulls
back to such a character of \(G\). Its Artin--Schreier break
is \(b_0\). Thus \(b_0\) is positive and prime to \(p\).
Finally
\[
 A-a=2p(p-1)>p(p-1)\ge p b_0.
\]
Lemma~\ref{lem:cancellation} applies. In particular the first
graded quotient has order \(p\), and with
\begin{equation}\label{eq:n-compositum}
 n=a+(p-1)b_0
\end{equation}
we have \(v_B(g\eta-\eta)=n+\ell_B(g)\) whenever
\(\ell_B(g)>b_0\).

\subsection{Exclude all early-conductor orderings}

Put \(F=F_2\), let \(b\) be its unique break, and write
\(u_2\) for the second upper break of \(L_2/K\).
Classical cyclic ramification gives \(u_2\ge pb\ge p\).
Let
\[
 E=FL_1,\qquad N=\Gal(B/E)=\langle\sigma^p\rangle .
\]
Then \(E/K\) is elementary abelian of degree \(p^2\).
The subgroup \(N\) persists in \(G^u\) for \(0\le u\le p-1\).
Indeed choose \(p-1<u'<u_2\).
The projection of \(G^{u'}\) on \(L_1\) is trivial, while its
projection on \(L_2\) contains \(\langle\sigma^p\rangle\).
In the direct product these two projections force all of
\(N\) to lie in \(G^{u'}\); monotonicity gives the assertion.
It follows that
\begin{equation}\label{eq:index-quotient}
 [G:G^u]=[G/N:(G/N)^u]\qquad(0\le u\le p-1).
\end{equation}
This is a quotient argument, not an assumed product rule for
compositum ramification.

Suppose \(b\le p-1\).
All degree-\(p\) characters of \(E/K\) are linear combinations
over \(\Fp\) of those for \(F/K\) and \(L_1/K\).
Their reduced polar representatives have largest poles at
most \(p-1\); a linear combination cannot introduce a larger
pole. The \(L_1\) character has break exactly \(p-1\).
Characters separate this elementary abelian group, so its
last upper break is \(p-1\).
By \eqref{eq:index-quotient} and the cancellation lemma its
first graded quotient has order \(p\). Hence \(E/K\) has
exactly two distinct upper breaks, \(b_0<p-1\) and \(p-1\),
each dropping the group by a factor \(p\).
This includes the possibility \(b=p-1\): possible
cancellation of the two leading poles either produces a
smaller first break, or gives a single drop of order \(p^2\),
which the lemma rules out.

The lower value corresponding to the last break of \(E/K\) is
\[
 c=\psi_B(p-1)=b_0+p(p-1-b_0)>b_0.
\]
There is \(h\in G\) with lower break exactly \(c\):
lift a nontrivial element of the last upper group of \(E/K\)
inside \(G^{p-1}\); its image disappears immediately after
that break. By \eqref{eq:n-compositum} and
\eqref{eq:later-rule},
\[
 v_B(h\eta-\eta)=n+c=a+p(p-1)<A,
\]
contradicting \eqref{eq:uniform-displ}. Therefore
\begin{equation}\label{eq:b-greater}
 b>p-1.
\end{equation}

\subsection{Read off the exact first quotient break}

By \eqref{eq:b-greater}, every nonzero linear combination
involving the character of \(F/K\) has break \(b\): its
highest pole cannot cancel with the strictly lower pole
of the \(L_1\) character. The other nonzero characters
have break \(p-1\). Thus \(E/K\) has upper breaks
\(p-1,b\), each of rank one.
The subgroup \(N\) persists through upper \(b\): choose
\(b<u'<u_2\), possible since \(u_2\ge pb>b\), and repeat
the two-projection argument.
Consequently
\[
 b_0=p-1,\qquad n=a+(p-1)^2.
\]
For \(p-1<u\le b\), the \(L_1\) projection is trivial
and the \(L_2\) projection is full. Hence
\(G^u=\langle\sigma\rangle\), and immediately after \(b\)
its \(L_2\) projection drops to \(\langle\sigma^p\rangle\).
The lower break of the native generator \(\sigma\) is thus
\[
 \ell_B(\sigma)=\psi_B(b)
 =(p-1)+p(b-(p-1))>b_0 .
\]
Its displacement is \(A\); the later-break rule yields
\[
 \ell_B(\sigma)=A-n
 =2p^2(p-1)-2p(p-1)^2-(p-1)^2=p^2-1.
\]
Comparison gives
\begin{equation}\label{eq:exact-first-break}
 (p-1)+p(b-(p-1))=p^2-1,\qquad b=2(p-1).
\end{equation}
Since \(F_e=F_2\) for \(e\ge2\), this proves the asserted
first quotient break at every higher level. Equivalently,
the largest pole of the native-oriented reduced AS class
is \(2(p-1)\), with nonzero coefficient. No general
formula for all its coefficients has been used or obtained.

\subsection{The complete second-layer filtration}

The full upper filtration of \(B/K\), at this point with
the last break \(u_2\) still undetermined, is
\begin{equation}\label{eq:compositum-filtration}
 G^u=
 \begin{cases}
 G,&0\le u\le p-1,\\
 \langle\sigma\rangle,&p-1<u\le 2(p-1),\\
 \langle\sigma^p\rangle,&2(p-1)<u\le u_2,\\
 1,&u>u_2.
 \end{cases}
\end{equation}
The first two lines were proved above. After \(p-1\) the
\(L_1\) projection is trivial, and upper quotient
compatibility identifies the \(L_2\) projection with its
upper group. A subgroup of \(C_{p^2}\times C_p\) with
trivial second projection is uniquely determined by its
first projection. This proves the last two lines without
any additional independence assumption.

Let \(c_0,c_1,c_2\) be the lower breaks of \(B/K\).
The first two are
\[
 c_0=p-1,\qquad c_1=(p-1)+p(p-1)=p^2-1.
\]
The element \(\tau=\sigma^p\) has lower break \(c_2\).
It fixes \(\beta\), so \eqref{eq:pstep} gives
\[
 v_B(\tau\eta-\eta)=2p^3(p-1).
\]
The later-break rule is applicable because \(c_2>c_1>b_0\).
Subtracting its identity for \(\sigma\) from that for \(\tau\)
gives
\[
 c_2-c_1=2p^3(p-1)-2p^2(p-1)=2p^2(p-1)^2.
\]
The index on the last interval of
\eqref{eq:compositum-filtration} is \(p^2\).
Hence Herbrand's formula gives
\[
 u_2-2(p-1)=\frac{c_2-c_1}{p^2}=2(p-1)^2,
 \qquad u_2=2p(p-1).
\]
In particular, the upper and lower breaks of \(B/K\) are
\begin{equation}\label{eq:compositum-breaks}
 \begin{aligned}
 &(p-1,\ 2(p-1),\ 2p(p-1)),\\
 &(p-1,\ p^2-1,\ p^2-1+2p^2(p-1)^2).
 \end{aligned}
\end{equation}
Upper quotient compatibility identifies the two upper
breaks of \(L_2/K\). Its lower breaks are
\[
 b_1=2(p-1),\qquad
 b_2=b_1+p\bigl(2p(p-1)-2(p-1)\bigr)
 =2(p-1)(p^2-p+1).
\]
These are actual invariants, derived with the actual first
break, independently of the hypothetical calculation in
Proposition~\ref{prop:separation}.

\subsection{Different exponents versus the root order}

For a finite totally ramified Galois extension \(D/K\), the
Hilbert different formula in the integer-normalized valuation is
\begin{equation}\label{eq:different-sum}
 d_{D/K}=\sum_{i\ge0}(|\Gal(D/K)_i|-1).
\end{equation}
One way to check this normalization is to use a uniformizer
\(\pi\), which generates the integer ring over \(R\).
Its minimal polynomial is Eisenstein; its derivative
generates the different. The valuation of that derivative
is \(\sum_{g\ne1}v_D(g\pi-\pi)\), and counting each
automorphism in its lower groups gives
\eqref{eq:different-sum}. The derivative description is
the monogenic case of the classical different formula
\cite[Tags 0BWD and 0BWG]{stacksdifferent}.

For \(L_2/K\), the successive lower group orders are
\(p^2,p\), giving
\[
 \begin{aligned}
 d_{L_2/K}
 &=(b_1+1)(p^2-1)+(b_2-b_1)(p-1)\\
 &=(2p-1)(p^2-1)+2p(p-1)^3\\
 &=(p-1)(2p^3-2p^2+3p-1).
 \end{aligned}
\]
This proves Theorem~\ref{thm:ramification}.
For \(B/K\), the three lower group orders are \(p^3,p^2,p\),
so the additional invariant is
\[
 \begin{aligned}
 d_{B/K}
 &=(c_0+1)(p^3-1)+(c_1-c_0)(p^2-1)
   +(c_2-c_1)(p-1)\\
 &=p(p^3-1)+p(p-1)(p^2-1)+2p^2(p-1)^3\\
 &=p^2(p-1)(2p^2-2p+3).
 \end{aligned}
\]

These are field different exponents, not the discriminant
of \(R[\alpha_2]\). For example, at \(p=3\) the proved
formulas give upper breaks \((4,12)\), lower breaks
\((4,28)\), and \(d_{L_2/K}=88\).
The root-difference product instead gives
\[
 v(\operatorname{Disc}(M_2))
 =p^2\bigl(p(p-1)2r+(p-1)2(p-1)\bigr)
 =4p^2(p-1)^2,
\]
which equals \(144\) at \(p=3\).
This is a direct consequence of the distances already
proved, not a numerical input to the field calculation.
It must not be substituted for the exponent \(88\).
